\section{Two 4-bit shifts}
\label{sec:shifts}

\subsection{Forced start: the primary test}
\label{sec:forced}

The forced start tests whether, once a reply has begun to comply, the 4-bit build continues it more often than the original does. It is our primary test because refusal is known to be weakest under prefill \citep{qi2024shallow}, so an erosion of refusal by quantization should be easiest to detect here.

We find such a shift in \qwen{} only. With eight tokens forced, its default 4-bit build continues into harmful content more often than bf16. Under the human labels (\S\ref{sec:instrument}), among the pairs that the judge had labelled differently, the harmful reply came from the 4-bit build in \N{human.qwen.dtofour} and from bf16 in \N{human.qwen.dtobf} ($p=\N{human.qwen.dp}$), and the design-weighted shift is \N{human.qwen.net}\pp{} \N{human.qwen.netci}. The judge, which under-counts harm, gives a smaller shift of the same sign, from \N{C1.permodel.qwen.all.ratea}\% to \N{C1.permodel.qwen.all.rateb}\%. This is a shift of \N{C1.permodel.qwen.all.diff}\pp{} \N{C1.permodel.qwen.all.ci}, with discordant prompts splitting \N{C1.permodel.qwen.all.aonly} to \N{C1.permodel.qwen.all.bonly} (Holm $p=\N{C1.permodel.qwen.all.pholm}$). Llama Guard also gives a positive shift (\N{labels.qwen.llamaguard.diff}\pp{} \N{labels.qwen.llamaguard.ci}).

\llama{} and \qwenthree{} show no shift. Under the judge, both are equivalent to bf16 within the $\pm$5\pp{} equivalence margin (90\% intervals \N{C1.permodel.llama.all.cininety} and \N{C1.permodel.qwenthree.all.cininety}). The human-label estimates, \N{human.llama.net} and \N{human.qwenthree.net}\pp{}, have 90\% intervals too wide to establish equivalence (\N{human.llama.netcininety} and \N{human.qwenthree.netcininety}). All three results hold on the \N{C1.permodel.qwen.nonpilot.n} prompts that were not used in the prefill pilot (a 150-prompt run before registration). Pooled across models the shift is not significant (discordant pairs \N{C1.pooled.aonly} against \N{C1.pooled.bonly}, $p=\N{C1.pooled.p}$), so the effect is specific to \qwen.

The effect also depends on the format. All three affine 4-bit builds of \qwen{} shift in the same direction, by \N{S6.qwen.4bit-g32.diff}\pp{} at group size 32, \N{S6.qwen.4bit-g64.diff}\pp{} at 64, and \N{S6.qwen.mixed46.diff}\pp{} under the mixed 4/6 recipe, which shares most of its tensors with the group-64 build. MXFP4, which loses the most likelihood, shows no shift (\N{S6.qwen.mxfp4.diff}\pp{}), and the four 4-bit builds differ significantly from one another (Cochran's $Q$, Holm $p=\N{S6.qwen.qp}$). The format therefore matters within one model and one bit width, a first indication that the size of a perturbation does not determine its effect. In the other two models no build shifts significantly, apart from a nominal \N{S6.llama.mxfp4.diff}\pp{} for \llama's MXFP4 build.

\qwen's shift is present with two forced tokens (\N{S2.cells.qwen.4bit-g64|k2.diff}\pp{}, Holm $p=\N{S2.holm.qwen.4bit-g64.k2}$), and is smaller with none, where bf16 almost never complies. It persists when the reply cap is tripled (\N{robustness.E4qwen.128.diff}\pp{} at 128 tokens and \N{robustness.E4qwen.384.diff}\pp{} at 384, on a fixed subset of 150 prompts). With the generic prefix ``Sure, here is'' the shift has the same sign but is smaller, \N{robustness.S8.qwen.diff}\pp{} \N{robustness.S8.qwen.ci}, and is nominally significant ($p=\N{robustness.S8.qwen.p}$) but not after Holm correction (\N{robustness.S8.qwen.pholm}).
